\documentclass[12pt,a4paper]{article}

\usepackage{fontspec}
\usepackage{unicode-math}
\usepackage[english]{babel}
\usepackage{geometry}
\usepackage{graphicx}
\usepackage{amsmath}
\usepackage{amssymb}
\usepackage{booktabs}
\usepackage{hyperref}
\usepackage{float}
\usepackage{caption}
\usepackage{subcaption}
\usepackage{xcolor}

\setmainfont{Libertinus Serif}
\setmathfont{Libertinus Math}

\geometry{left=2.5cm,right=2.5cm,top=2.5cm,bottom=2.5cm}

\hypersetup{
    colorlinks=true,
    linkcolor=blue!50!black,
    urlcolor=blue!50!black,
    citecolor=blue!50!black,
    pdftitle={Lab 10 Report - HOG and Applications},
    pdfauthor={Bui Quang Chien}
}

\setcounter{tocdepth}{1}
\setcounter{secnumdepth}{3}

\newcommand{\labfig}[3]{%
\begin{figure}[H]
    \centering
    \IfFileExists{#1}{%
        \includegraphics[width=0.95\linewidth]{#1}%
    }{%
        \fbox{\parbox{0.9\linewidth}{Run notebook export cell to generate #1}}%
    }
    \caption{#2}
    \label{#3}
\end{figure}
}

\begin{document}

\begin{titlepage}
    \centering
    \vspace*{2cm}

    {\LARGE\bfseries HANOI UNIVERSITY OF SCIENCE\\[0.3cm]
    FACULTY OF MATHEMATICS, MECHANICS AND INFORMATICS\\[0.8cm]}

    {\huge\bfseries LAB REPORT\\[0.4cm]
    COMPUTER VISION (MAT3562E 2)\\[0.4cm]}

    {\Large\bfseries LAB 10: HISTOGRAM OF ORIENTED GRADIENTS (HOG)\\
    AND ITS APPLICATIONS\\[1.2cm]}

    \begin{flushleft}
    {\large
    Lecturer: Dr. Cao Van Chung\\
    Student: Bui Quang Chien\\
    Student ID: 23001837\\
    Class: MAT3562E 2
    }
    \end{flushleft}

    \vfill
    {\large Hanoi, April 2026}
\end{titlepage}

\tableofcontents
\newpage

\section{Overview and Learning Objectives}
This report completes all core requirements for Lab 10:
\begin{enumerate}
    \item Manual HOG implementation from scratch using NumPy.
    \item Comparison against \texttt{scikit-image} HOG.
    \item Two practical applications:
    image classification and classical object detection.
    \item Error analysis with quantitative evidence.
\end{enumerate}

The implementation is in:
\texttt{lab10\_BuiQuangChien\_23001837.ipynb}.

\section{Dataset and Experimental Setup}
\subsection{Data used}
\begin{itemize}
    \item Traffic sign image (and fallback camera image) for manual HOG visualization.
    \item \texttt{sklearn.datasets.load\_digits()} for quantitative classification and detector experiments.
\end{itemize}

\subsection{HOG configuration}
The same HOG configuration is used across manual and library pipelines:
\begin{itemize}
    \item Orientation bins: 9.
    \item Cell size: $8\times8$.
    \item Block size: $2\times2$ cells.
    \item Block normalization: L2-Hys.
\end{itemize}

\subsection{Core equations}
Given Sobel gradients $G_x, G_y$:
\begin{align}
M(x,y) &= \sqrt{G_x(x,y)^2 + G_y(x,y)^2},\\
\theta(x,y) &= \left(\operatorname{atan2}(G_y, G_x)\cdot \frac{180}{\pi} + 180\right) \bmod 180.
\end{align}
For each block vector $v$, L2-Hys normalization is:
\begin{align}
\hat{v} &= \frac{v}{\sqrt{\|v\|_2^2 + \varepsilon^2}},\\
\tilde{v} &= \min(\hat{v}, 0.2),\\
v_{\text{L2-Hys}} &= \frac{\tilde{v}}{\sqrt{\|\tilde{v}\|_2^2 + \varepsilon^2}}.
\end{align}

\section{Manual HOG Implementation}
The manual pipeline includes:
\begin{enumerate}
    \item Preprocessing and resize.
    \item Sobel gradient computation.
    \item Orientation quantization and magnitude voting.
    \item Histogram per cell.
    \item Block normalization and feature concatenation.
    \item Visualization of dominant local orientations.
\end{enumerate}

\labfig{outputs/manual_hog_pipeline.png}{Manual HOG pipeline visualization (input, gradients, magnitude, orientation, HOG view).}{fig:manual_hog}

On the $128\times128$ sample image, the manual feature vector length is 8100, matching library output.

\section{Manual vs Library Comparison}
\subsection{Feature dimensions and speed}
\begin{table}[H]
\centering
\caption{Manual HOG vs \texttt{scikit-image} HOG}
\begin{tabular}{lcc}
\toprule
Metric & Manual & \texttt{scikit-image} \\
\midrule
Feature length (sample image) & 8100 & 8100 \\
Single-image extraction time (ms) & 10.80 & 4.19 \\
Batch extraction time (ms/image) & 0.790 & 0.190 \\
\bottomrule
\end{tabular}
\end{table}

The library implementation is faster, while the manual implementation provides full algorithmic transparency.

\section{Application 1: Image Classification}
\subsection{Pipeline}
\begin{enumerate}
    \item Upscale digits to $32\times32$.
    \item Extract HOG features (manual or library).
    \item Train classifiers: LinearSVC and Logistic Regression.
    \item Evaluate using accuracy and confusion matrix.
\end{enumerate}

\subsection{Model comparison}
\begin{table}[H]
\centering
\caption{Classification results from \texttt{outputs/classification\_model\_comparison.csv}}
\begin{tabular}{lccc}
\toprule
Feature + Classifier & Accuracy & Train time (s) & Feature dim \\
\midrule
Manual + LogisticRegression & 0.9889 & 0.5568 & 324 \\
Skimage + LogisticRegression & 0.9889 & 0.2579 & 324 \\
Manual + LinearSVC & 0.9822 & 4.0223 & 324 \\
Skimage + LinearSVC & 0.9756 & 2.9581 & 324 \\
\bottomrule
\end{tabular}
\end{table}

Best classification configuration: Manual HOG + Logistic Regression with accuracy $0.9889$.

\labfig{outputs/classification_confusion_matrix.png}{Confusion matrix of the best classifier.}{fig:cm}

\section{Application 2: Classical HOG Object Detector}
\subsection{Detector design}
The detector follows the classical pipeline:
\begin{enumerate}
    \item Image pyramid for multi-scale search.
    \item Sliding window scan over each scale.
    \item HOG extraction and linear SVM scoring.
    \item Non-Maximum Suppression (NMS) with IoU thresholding.
\end{enumerate}

\subsection{Detector classifier quality}
On the binary detector training split (digit 0 vs non-0):
\begin{itemize}
    \item Accuracy: 0.9900
    \item Precision: 0.9706
    \item Recall: 1.0000
    \item F1-score: 0.9851
    \item PR-AUC: 0.9991
\end{itemize}

\subsection{Threshold tuning to reduce false positives}
Grid search over detector thresholds on 10 synthetic scenes gives best configuration:
\begin{itemize}
    \item Score threshold: 0.75
    \item NMS IoU threshold: 0.15
    \item Aggregated precision: 0.3548
    \item Aggregated recall: 0.8800
    \item Aggregated F1: 0.5057
\end{itemize}

\begin{table}[H]
\centering
\caption{Single-scene comparison: baseline vs tuned thresholds}
\begin{tabular}{lcccccc}
\toprule
Setting & TP & FP & FN & Precision & Recall & F1 \\
\midrule
Baseline $(0.35, 0.25)$ & 5 & 15 & 0 & 0.250 & 1.000 & 0.400 \\
Tuned $(0.75, 0.15)$ & 5 & 7 & 0 & 0.417 & 1.000 & 0.588 \\
\bottomrule
\end{tabular}
\end{table}

\labfig{outputs/detector_baseline_vs_tuned.png}{Detector output: baseline thresholds vs tuned thresholds.}{fig:detector_compare}

\section{Error Analysis}
\subsection{Classification errors}
Per-class metrics remain high for all classes (lowest F1 around 0.9773).
Top confusion pairs are sparse and each appears once:
\begin{itemize}
    \item $3 \rightarrow 1$
    \item $5 \rightarrow 9$
    \item $6 \rightarrow 1$
    \item $6 \rightarrow 8$
    \item $9 \rightarrow 8$
\end{itemize}
This indicates strong class separation from HOG features on the digits task.

\subsection{Detection errors}
From \texttt{outputs/detection\_summary\_metrics.csv}:
\begin{table}[H]
\centering
\caption{Aggregated detection error summary after tuning}
\begin{tabular}{lc}
\toprule
Metric & Value \\
\midrule
Total TP & 44 \\
Total FP (near-miss) & 6 \\
Total FP (background) & 74 \\
Total FN & 6 \\
Precision & 0.3548 \\
Recall & 0.8800 \\
F1 & 0.5057 \\
Mean IoU on TP & 0.6669 \\
FN small / medium / large & 3 / 3 / 0 \\
\bottomrule
\end{tabular}
\end{table}

\labfig{outputs/detector_error_breakdown.png}{False-positive composition and false-negative size distribution.}{fig:error_breakdown}

Main observations:
\begin{itemize}
    \item Background false positives dominate near-miss false positives.
    \item False negatives mostly appear in small and medium objects; no large-object misses.
    \item This is consistent with HOG sensitivity to local edge patterns and sliding-window discretization.
\end{itemize}

\section{Conclusion}
Manual HOG implementation clarified every algorithmic stage and matched the library feature length.
The library version is computationally faster.
For classification, HOG + linear models achieved very high accuracy.
For classical detection, tuning score and NMS thresholds reduced duplicate detections and improved precision on the inspected scene while keeping recall high.

\section{Reproducibility Artifacts}
Generated files under \texttt{lab-10/outputs}:
\begin{itemize}
    \item PNG: \texttt{manual\_hog\_pipeline.png}, \texttt{classification\_confusion\_matrix.png}, \texttt{detector\_baseline\_vs\_tuned.png}, \texttt{detector\_error\_breakdown.png}
    \item CSV: \texttt{classification\_model\_comparison.csv}, \texttt{detector\_threshold\_tuning.csv}, \texttt{classification\_per\_class\_metrics.csv}, \texttt{classification\_top\_confusions.csv}, \texttt{detection\_scene\_level\_breakdown.csv}, \texttt{detection\_summary\_metrics.csv}
\end{itemize}

\end{document}
